El objetivo general (Sección~\ref{secc:og}) se cumplió.
El sistema se construyó y se midió contra decisores simples, contra una sola memoria con
las ocho clases y contra un codificador sin la pérdida de clase.
Lo que dejaron esas comparaciones se resume por objetivo específico.

\paragraph{Arquitectura (objetivo~\ref{oe:arquitectura}).}
Son ocho módulos, cada uno con una memoria de contenido heteroasociativa y dos
directorios, y un mediador que deja de intervenir en la fase madura.
Registrar, reconocer, recuperar y enrutar son operaciones de las memorias
(Sección~\ref{sec:operaciones}); la suma por tokens, la lectura B1, el umbral $\tau$ y el
rescate léxico operan fuera de ellas (Sección~\ref{sec:auditoria}).
Una sola memoria de contenido por módulo basta: los pesos de cada dominio se leen de la
memoria heteroasociativa, y con ellos el sistema toma las mismas decisiones que con
memorias homoasociativas separadas en imagen y casi todas en texto (cambian algunos
destinos del protocolo y una consulta de 80 en la curva de capacidad;
Sección~\ref{sec:una-memoria}).

\paragraph{Formación del directorio (objetivo~\ref{oe:formacion}).}
Sobre conjuntos de 120 consultas que nunca forman el directorio, la fase madura con la
lectura B1 sube del 78.9\,\% con 50 consultas de formación al 91.0\,\% con 291, y la
temprana es del 89.7\,\% (promedio de cinco sorteos del conjunto de evaluación).
Con unas 200 consultas, la fase madura alcanza a la temprana: la diferencia es de
$+0.2$ puntos con 200 y de $+1.3$ con 291, y ninguna diferencia positiva se sostiene
entre sorteos.
El directorio coincide, entonces, con la decisión del mediador en el 92.5\,\% de las
consultas evaluadas (sorteo principal), con un acierto equivalente y sin difundir la
consulta a todos los módulos.
La lectura B1 supera a la cruda con todos los tamaños de formación y en todos los
sorteos, por 4.5 a 8.8 puntos en promedio; sin ella, la fase madura nunca alcanza a la
temprana.

\paragraph{Frente a decisores simples y a una memoria única
(objetivo~\ref{oe:lineas-base}).}
En texto, el sistema acierta lo mismo que los decisores simples.
Sobre las 80 consultas del protocolo completo, en las que solo 11 traen el nombre de
su clase, acierta el 81.2\,\%, y el voto por tokens del vocabulario, el centroide y los
vecinos más cercanos dan entre 76.0 y 81.5\,\%.
Sobre las 171 consultas reservadas, con 200 imágenes por clase, enruta bien el
96.8\,\% y la memoria única clasifica bien el 96.9\,\% (diferencia de $-0.1$ puntos,
$p = 0.86$).
Lo que la memoria agrega sobre el voto en texto es el desempate entre las clases que
contienen algún token de la consulta: 2.9 puntos sobre las 400 consultas de
calibración (90.0 frente a 87.1\,\%).
Una tabla de conteos de tokens formada con las mismas consultas no se distingue del
directorio (92.2 frente a 93.0\,\%, $p = 0.58$).

\paragraph{Rechazo (objetivo~\ref{oe:rechazo}).}
De las 40 consultas fuera de dominio, el sistema, con la doble compuerta, acepta 2 con
25, 50, 100 y 200 imágenes por clase, y la memoria única 2, 2, 5 y 8; la segunda
compuerta se decidió después de ver estas consultas (limitación (g)).
Los decisores de coseno aceptan las 40.
La compuerta textual coincide con la pertenencia al vocabulario en los 3064 pares de
token y clase examinados, de modo que un voto léxico rechaza las mismas consultas.
Las dos que pasan son las únicas que comparten una palabra con el vocabulario de un
módulo (limitación (a)).
En imagen, el umbral de energía y la contención rechazan buena parte de las entradas sin
estructura (limitación (c)).
El rechazo de imágenes reales de otras categorías no se midió.

\paragraph{Recuperación en ambas direcciones (objetivo~\ref{oe:recuperacion}).}
De texto a imagen, el sistema responde con una imagen al 98.7\,\% de las 171 consultas
reservadas, y la memoria única al 100\,\%.
La distancia de la imagen recordada a la imagen real más cercana de su clase es de
27.2 con 200 imágenes por clase, frente a 32.5 de la memoria única y 26.8 del oráculo.
De imagen a texto, 113 de las 116 imágenes enrutadas entre las 15 primeras de cada clase
evocan palabras de su dominio (97.4\,\%).

\paragraph{Imágenes y codificador (objetivo~\ref{oe:vision}).}
Los directorios visuales envían al especialista el 96.2\,\% de las 656 imágenes de
prueba (631), rechazan 23 y se equivocan en dos.
El centroide y los vecinos más cercanos sobre los mismos códigos aciertan de 652 a 655
sin rechazar ninguna, y el vecino más cercano acierta las 23 que el sistema rechaza.
Con un codificador entrenado sin la pérdida de clase, el enrutamiento cae al 70.4\,\%,
los falsos enrutamientos suben de 2 a 169 y la cobertura de punta a punta con objetos nuevos baja del
89.2 al 59.3\,\%.
La separación entre clases la hace el codificador supervisado, y la propiedad de
fallar con un rechazo en lugar de un error depende de esa supervisión.
Con objetos que ningún módulo vio, la cobertura es del 89.2\,\% frente al 95.5\,\% con
vistas nuevas de objetos conocidos, pero el codificador sí había visto esos objetos;
con un codificador reentrenado sin el objeto reservado, la cobertura es del 87.6\,\%
[84.3, 91.0] y los falsos enrutamientos suben de 2 a 136 de 3280 (4.1\,\%).

\paragraph{En suma.}
El sistema enruta como un decisor simple y, además, rechaza lo que no contiene,
recupera en ambas direcciones y guarda un registro de quién sabe qué.
Su acierto lo ponen las representaciones: el vocabulario en texto y el codificador
supervisado en imagen.
